\section{The complete small quadratic family at coefficient seven}
\label{sec:coefficient-seven-full-quadratic-window}

\paragraph{The exact preceding remainder.}
Retain the complete original datum and remainder of
AC-G200-window, including every coefficient-five omission.
In particular $h\geq2$ and $T\geq1$ are integers, and on an
actual source page the parameters are
\[
 \begin{gathered}
 H_u=\mathbb Q(\zeta_{2^{h+2}},\zeta_7,\sqrt{-q},\sqrt{-d},
                 \omega,\sqrt[3]{269}),\qquad K=\mathbb Q(\omega),\\
 u^{2+8\varepsilon}\leq2^h\leq2u^{2+8\varepsilon},\qquad
 T=\lfloor u^{399}\rfloor,\qquad0<\varepsilon<1/256.
 \end{gathered}\tag{S7A.1}
\]
Keep exactly $Z_u,J_{u,0},J_{u,\rm mix},V_u,Y_u,C_7(u),U_7(u)$
and $C_5(u),Z_5(u),U_5(u)$ from (F5W.2)--(F5W.6) and their
cited preceding definitions. In particular $Y_7(u)\subset Y_u$
is the single coefficient-seven family
$(7,\tau_7^e|_{G_K},e)$, one triple per weight.
The rings $O_\ell$, fraction fields $E_\ell$, normalizations
$b_\ell$, lattices and every full local condition remain the
ones of (FC.1)--(FC.2).

Put $\mathcal D=\{1,\chi_{-4},\chi_8,\chi_{-8}\}$ and use
$\tau_7$ with its specified original values in $\mathbb Z_7$.
Define actual triples, with every undisplayed cubic and auxiliary
factor trivial, by
\[
 \begin{split}
 B_7(u)=\{\,&(7,(\chi_8\tau_7^j)|_{G_K},e):
              1\leq e<T,\ 0\leq j<6,\ j-e\text{ even}\},\\
 C_7^{\rm all}(u)=\{\,&(7,(\delta\tau_7^j)|_{G_K},e):
          1\leq e<T,\ \delta\in\mathcal D,\ 0\leq j<6\},\\
 U_7^{\rm all}(u)&=C_7^{\rm all}(u)\setminus Y_7(u).
 \end{split}\tag{S7A.2}
\]
The word ``all'' refers only to these four two-power factors and
six seven-parts, with trivial cubic and auxiliary factors.
It does not refer to every character at coefficient seven.

For any actual trivial-cubic character $\eta$, let
$\psi_{\eta,e}^{\rm nc}$ be its unique rational finite extension
with sign $(-1)^e$, using the same original coefficient.
Define the explicit nonnegative remainder
\[
 \begin{split}
 \mathcal R_{399}^{\rm match,7all,5}(u)={}&
 \sum_{(\ell,\eta,e)\in J_{u,\rm mix}}
                  \frac{\log\#\mathcal S_{\ell,\eta}(e)}{b_\ell}\\
 &+\sum_{(\ell,\eta,e)\in
    J_{u,0}\setminus(V_u\cup Y_u\cup C_7^{\rm all}(u)\cup C_5(u))}
       \frac{\log\#\mathcal S^{\mathbb Q}_{\ell,
                                  \psi_{\eta,e}^{\rm nc}}(e)}{b_\ell}.
 \end{split}\tag{S7A.3}
\]
All its groups are finite by the retained AC-G191/AC-G192 proof.

\begin{theorem}[AC-G202-window: completing the small coefficient-seven family]
\label{AC-G202-window}
The set $B_7(u)$ consists of actual original-host triples and
\[
 \begin{gathered}
 C_7^{\rm all}(u)=C_7(u)\sqcup B_7(u),\qquad
 U_7^{\rm all}(u)=U_7(u)\sqcup B_7(u),\\
 B_7(u)\cap\bigl(Z_u\cup V_u\cup Y_u\cup C_7(u)\cup C_5(u)\bigr)
                      =\varnothing,\\
 \mathcal R_{399}^{\rm match,7all,5}(u)
                        =\mathcal R_{399}^{\rm match,7,5}(u).
 \end{gathered}\tag{S7A.4}
\]
The equality deletes zero rational noncritical terms. It does not
assert zero for their original-$K$ strict groups or critical partners.
For $N=T-1$ the exact counts are
\[
 \begin{aligned}
 |B_7(u)|&=3N,& |C_7^{\rm all}(u)|&=24N,\\
 |U_7^{\rm all}(u)|&=23N,& |U_5(u)|&=24N-10\lfloor N/4\rfloor.
 \end{aligned}\tag{S7A.5}
\]
Thus exactly $3N$ further rational terms disappear relative to
(F5W.5), and $27N-10\lfloor N/4\rfloor$ relative to (S7W.5).
The original triples no longer represented form the disjoint union
\[
 Z_u\sqcup V_u\sqcup Y_u\sqcup U_7^{\rm all}(u)\sqcup U_5(u),
 \tag{S7A.6}
\]
whose cardinality is
$|Z_u|+2\mathbf1_{h\geq3}N+|Y_u|+47N-10\lfloor N/4\rfloor$.
For $T\geq2$ the indexing-set reduction is proper; there is no
claim of strict numerical decrease between equal remainders.

Retain exactly the costs and constants of (F5W.8). Then
\[
 \begin{aligned}
 |I_{399}-\mathcal R_{399}^{\rm match,7all,5}|&\leq
      \mathcal E_u+\mathcal L_u+\mathcal A_u+C_2\mathcal P_u,\\
 |I_{399}^{\rm all7,2}-\mathcal R_{399}^{\rm match,7all,5}|&\leq
      \mathcal E_u+\mathcal L_u+C_2\mathcal P_u.
 \end{aligned}\tag{S7A.7}
\]
One full (FL.4) majorant, one (LT.2) torsion sum and one critical
aggregate still suffice. No new family multiplier or error occurs.
The envelope remains
$O_\varepsilon(u^{806+16\varepsilon}(\log u)^2)$.
The old and new numerical remainder premises have the identical
window, threshold and slack. Their sufficient comparison with
the full original index premise retains the same threshold
increase and half-slack construction.
\end{theorem}

\begin{proof}
\proofstep{1. Actual characters and the original noncritical coefficient.}
The original host contains $\sqrt2$ and $\mathbb Q(\zeta_7)$.
The six possible values of $\tau_7^j$ already lie in the original
$O_7$, and each character order divides $3\cdot2^h$.
No coefficient extension, alternative residue embedding or new
finite character is introduced. Distinct two-parts and seven-parts
remain distinct over $K$: their rational ratio has conductor
prime to three and cannot be $\eta_K=\chi_{-3}$.

For a triple in $B_7(u)$ the finite character
$\chi_8\tau_7^j$ has sign $(-1)^j=(-1)^e$.
It is therefore already the exact noncritical rational extension,
without an extra factor of $\eta_K$.
The original integral theorem (RP7.5) gives
\[
       \mathcal S^{\mathbb Q}_7\bigl((E_7/O_7)
                       (\chi_8\tau_7^j c_7^{-e})\bigr)=0.
 \tag{S7A.8}
\]
Its proof uses the actual real field
$\mathbb Q(\zeta_7)^+\mathbb Q(\sqrt2)$, its two seven-adic
places and surjectivity of each individual local group onto the
maximal pro-seven quotient. It retains one full local coboundary
witness killed by an original uniformizer power.
Restriction--corestriction uses degree six, a unit in $O_7$.
No one-prime assertion is substituted for that two-prime proof.

The full index-two original lattice comparison in (RP7.5)
keeps both completions of $K$ at seven, all away-inertia
conditions and the real convention. It identifies the original
$K$ strict group with its critical rational partner, not with zero.
That partner already belongs to the single aggregate in (CP.2).
The same finite-localization inflation and finite-etale coefficient
dictionary of (FC.2) remain in force.

\proofstep{2. Exhaustive old overlap and exact counts.}
At coefficient seven the old $Z_u$ has only the primitive high
two-power family from (OPW.3), whose threshold is $c_7=16$.
Indeed $v_2(7^2-1)=4$ and
$2^{15}\leq21\cdot7\cdot2^8=37632<2^{16}$.
Every new two-part has conductor eight, so those families
are disjoint. The old inert-trivial family does not contain
seven, which is split modulo three. All other old $Z_u$
coefficient families and $V_u$ have different coefficients.

Among the families of $Y_u$, the only coefficient-seven one
is $Y_7(u)$, whose two-part is trivial. It is disjoint from
$B_7(u)$, and is contained in $C_7(u)$ through the $\delta=1$
rows. The coefficient-five set $C_5(u)$ is disjoint by its
coefficient. The preceding $C_7(u)$ includes all six seven-parts
for $\delta=1,\chi_{-4},\chi_{-8}$, but only the three rows
with $j-e$ odd for $\delta=\chi_8$. Its missing rows are exactly
$B_7(u)$. This proves every set identity in (S7A.4).

There are three $j$ of each parity at every weight.
Thus $|B_7|=3N$, the full family has $4\cdot6N=24N$,
and deleting the one old $Y_7$ row per weight gives $23N$.
Retaining the exact G200 count for $U_5$ proves (S7A.5) and
(S7A.6), including $T=1$. At $T\geq2$, the triple
$(7,(\chi_8\tau_7)|_{G_K},1)$ is a new row, proving properness.

\proofstep{3. Zero-term deletion and the one existing cost.}
Apply (S7A.8) to each newly absent term. Every original
normalization $b_7$ is the same, and all sums are finite.
Equation (RP7.7) gives exactly the last line of (S7A.4).
No original positive derived index has been separated into
rational positive indices, and no critical partner is discarded.
Substituting this pointwise equality into (F5W.8) proves
(S7A.7) with the identical constants.
The ambient degree $D_u=144\cdot2^h$ and abelian degree
$d_u=96\cdot2^h$ keep their separate roles in those costs.

\proofstep{4. The identical common window.}
Fix the original datum, page and $0<\varepsilon<1/256$.
Retain the (TC.3) parameters $J,X_0,s$, every prescribed pair,
and the same window $X^{1/J}\leq W$, $W^2\leq X$.
The equality (S7A.4) makes the two explicit numerical remainder
premises equivalent with exactly the same threshold and slack.
For comparison to the full original index use the same uniform
$U_0$ making the first error in (S7A.7) at most
$(s/2)\exp(u^\varepsilon/2)$, the same raised threshold
$\max(X_0,U_0^J)$ and the same $W$ with slack $s/2$.
No independent existential windows are intersected.
All actual source-rank $R.t$, field, coefficient, lattice,
page, pair and simultaneous finite constraints are unchanged.
\end{proof}

\paragraph{Remaining scope.}
The nine unsupported coefficient-five rows at weights divisible
by four and every other remaining mixed or noncritical term are
still present. The explicit numerical remainder estimate,
mixed theta/Selmer realization and H04.4 remain open.
No further claim is made for other coefficient-seven characters.
There is no new question row or independent helper-credit claim;
completed reductions are assessed together at their actual times.
